\documentclass[11pt]{amsart}
\usepackage[a4paper,margin=1in]{geometry}
\usepackage{amsmath,amssymb,amsthm,mathtools}
\usepackage[T1]{fontenc}
\usepackage{lmodern}
\usepackage{microtype}
\usepackage{xcolor}
\usepackage{hyperref}
\hypersetup{colorlinks=true,linkcolor=blue!50!black,citecolor=blue!50!black,urlcolor=blue!50!black}

\newtheorem{theorem}{Theorem}

\newcommand{\E}{\mathbb{E}}
\newcommand{\optprior}{\mathrm{OPT}_{\mathrm{prior}}}
\newcommand{\optadapt}{\mathrm{OPT}_{\mathrm{adapt}}}

\title{A $5/4$ Lower-Bound Family for the A Priori / Adaptive Gap}
\author{}
\date{}

\begin{document}
\maketitle

We work with the same single-cycle subdivision family as in \texttt{single\_cycle\_ratio.tex}. Let $n\ge 7$ be odd, let $v_1,\dots,v_n$ be the vertices of $K_n$ (indices modulo $n$), and fix the Hamiltonian cycle
\[
C:=v_1v_2\cdots v_nv_1.
\]
For each $i\in[n]$, subdivide the edge $v_iv_{i+1}$ by a new vertex $x_i$, and let $d$ be the shortest-path metric of the resulting graph. We take $r:=v_1$ as depot. The client set is
\[
V:=\{v_2,\dots,v_n,x_1,\dots,x_n\},
\]
so the depot is $r=v_1\notin V$. Every original client $v_2,\dots,v_n$ is active with probability $1$, and each subdividing vertex $x_i$ is active independently with probability $1/2$.

The goal of this note is to prove that this family also yields the asymptotic lower bound $5/4$ for the ratio between the optimal a priori and adaptive values.

\begin{theorem}\label{thm:main}
For every odd $n\ge 7$, the above single-cycle subdivision instance satisfies
\[
\optadapt\le \frac32 n+18.
\]
Consequently, using the previously established identity
\[
\optprior=\frac{15}{8}n,
\]
one has
\[
\frac{\optprior}{\optadapt}
\ge
\frac{(15/8)n}{(3/2)n+18}
\xrightarrow[n\to\infty]{}\frac54.
\]
In particular, for every $\varepsilon>0$, for all sufficiently large odd $n$,
\[
\frac{\optprior}{\optadapt}\ge \frac54-\varepsilon.
\]
\end{theorem}

\begin{proof}
We define an explicit adaptive policy $\pi$ and bound its expected cost.

Let
\[
M:=\{v_3,\dots,v_{n-1},x_3,\dots,x_{n-2}\}
\]
be the set of \emph{middle} clients, and let
\[
B:=V\setminus M=\{v_2,v_n,x_1,x_2,x_{n-1},x_n\}
\]
be the set of \emph{boundary} clients.

\medskip
\noindent\textbf{Step 1: visit all boundary clients first.}
The policy first calls all six vertices of $B$ in an arbitrary fixed order, and then returns to the depot $r$.

Every distance in this metric is at most $3$, so this boundary phase has cost at most
\[
7\cdot 3=21,
\]
independently of the realization.

\medskip
\noindent\textbf{Step 2: the alternating middle policy.}
After finishing the boundary phase, the policy is back at $r$. It then calls $v_3$, moves to $v_3$, and from that point on works only on the middle path
\[
P:=v_3,x_3,v_4,\dots,x_{n-2},v_{n-1}.
\]

At every stage, the uncalled middle clients form a subpath of $P$. The current position is one of the two endpoint originals of this remaining subpath. If the current position is the left endpoint original, the policy probes rightward: it calls the adjacent subdividing vertex. If that subdividing vertex is active, then the policy immediately calls the next original vertex and continues in the same direction. If that subdividing vertex is inactive, then the policy jumps to the right endpoint original of the remaining subpath and proceeds symmetrically leftward. The same rule is applied when the current position is the right endpoint original.

Equivalently, the policy keeps moving along the current side of the remaining subpath for as long as the encountered middle subdividing vertices are active, and each time it meets an inactive middle subdividing vertex it switches to the opposite side of the remaining subpath.

\medskip
\noindent\textbf{Step 3: deterministic cost of the middle phase.}
Fix a realization $A$, and let
\[
a_M:=|A\cap\{x_3,\dots,x_{n-2}\}|
\]
be the number of active middle subdividing vertices. Since there are exactly $n-4$ middle subdividing vertices, the number of inactive middle subdividing vertices is
\[
(n-4)-a_M.
\]

We claim that the total cost of the middle phase, including the initial move from $r$ to $v_3$ and the final return to $r$, is at most
\[
n+a_M-1.
\]
Indeed:

\begin{itemize}
    \item The initial move from $r=v_1$ to $v_3$ costs $1$.
    \item Each active middle subdividing vertex contributes exactly $2$: one unit to move from an endpoint original to that subdividing vertex, and one unit to move from there to the next original. Hence all active middle subdividing vertices together contribute exactly
    \[
    2a_M.
    \]
    \item Each inactive middle subdividing vertex causes one switch from one end of the remaining middle subpath to the other. Such a switch costs $1$, except that the very last switch may cost $2$ when only one middle original vertex remains uncalled. Therefore all switches together cost at most
    \[
    ((n-4)-a_M)+1.
    \]
    \item After all middle clients have been called, the salesperson is at some original vertex, and returning from any original vertex to $r$ costs at most $1$.
\end{itemize}

Summing these contributions gives
\[
1+2a_M+\bigl((n-4)-a_M+1\bigr)+1=n+a_M-1.
\]
This proves the claim.

\medskip
\noindent\textbf{Step 4: expectation.}
For every realization $A$, the total cost of $\pi$ is at most
\[
21+(n+a_M-1)=n+a_M+20.
\]
Taking expectation and using $\E[a_M]=(n-4)/2$, we obtain
\[
\E[c(T(A;\pi))]
\le n+\frac{n-4}{2}+20
=\frac32 n+18.
\]
Hence
\[
\optadapt\le \frac32 n+18.
\]
The ratio statement follows immediately from the previously proved identity $\optprior=\frac{15}{8}n$.
\end{proof}

\end{document}
